\section{编写 Skills 的最佳实践}

确定了要做什么类型的 Skill 之后，接下来的问题是\textbf{怎么写}。以下是 Anthropic 总结的最佳实践、技巧和窍门。Anthropic 最近还发布了 Skill Creator 工具，帮助在 Claude Code 中更方便地创建 Skills。

\subsection{不要说显而易见的事（Don't State the Obvious）}

Claude Code 对代码库已经非常了解，Claude 本身对编程也很在行，有很多默认的观点和偏好。如果你要发布一个主要提供知识的 Skill，\textbf{应当聚焦于那些能把 Claude 推离其常规思维方式的信息}。

\begin{knowledgebox}{frontend-design Skill 案例}
这个 Skill 由 Anthropic 的一位工程师开发，通过与客户反复迭代来改进 Claude 的设计品味。它专门帮 Claude 避免那些经典的套路模式（classic patterns），比如 Inter 字体和紫色渐变。这是一个非常好的``非显而易见知识'' Skill 的范例：它提供的不是 Claude 已经知道的东西，而是纠正 Claude 的默认倾向。
\end{knowledgebox}

\subsection{建一个踩坑点章节（Build a Gotchas Section）}

\begin{importantbox}{Gotchas 是 Skill 中信息量最高的部分}
任何 Skill 中最有价值的内容就是踩坑点（Gotchas）章节。这些章节应当根据 Claude 在使用你的 Skill 时遇到的常见失败点\textbf{逐步积累}。理想情况下，你会随着时间推移持续更新 Skill，把新发现的 gotchas 补充进去。
\end{importantbox}

\subsection{利用文件系统与渐进式披露（Use the File System \& Progressive Disclosure）}

Skill 是一个文件夹，不只是一个 markdown 文件。你应该把\textbf{整个文件系统}当作 context engineering（上下文工程）和 progressive disclosure（渐进式披露）的工具。告诉 Claude 你的 Skill 里有哪些文件，它会在合适的时候去读取它们。

\begin{knowledgebox}{Context Engineering 与 Progressive Disclosure}
\textbf{Context engineering} 是指精心设计和管理输入给大语言模型的上下文信息，以最大化输出质量。\textbf{Progressive disclosure} 借用 UI 设计概念，指不一次性把所有信息塞给模型，而是让它在需要时再去读取，从而节省上下文窗口空间并减少干扰。
\end{knowledgebox}

渐进式披露的具体形式：
\begin{itemize}
    \item \textbf{最简形式}：指向其他 markdown 文件让 Claude 按需读取。例如，把详细的函数签名和用法示例拆分到 \texttt{references/api.md}
    \item \textbf{模板文件}：如果最终输出是一个 markdown 文件，可以在 \texttt{assets/} 中放一个模板文件供复制使用
    \item \textbf{资源文件夹}：提供 references、scripts、examples 等文件夹，帮助 Claude 更高效地工作
\end{itemize}

\subsection{不要把 Claude 限制得太死（Avoid Railroading Claude）}

\begin{warningbox}{过度约束的指令会降低 Skill 的复用性}
Claude 通常会努力严格遵循你的指令，而 Skills 的复用性很强，所以你需要\textbf{小心不要在指令中过于具体}。给 Claude 它需要的信息，但留给它适应具体情况的灵活性（flexibility to adapt to the situation）。
\end{warningbox}

\subsection{考虑好初始设置（Think through the Setup）}

有些 Skills 需要用户提供上下文来完成初始设置。例如，一个发站会汇报到 Slack 的 Skill 可能需要知道发到哪个频道。

推荐做法：
\begin{itemize}
    \item 把设置信息存在 Skill 目录下的 \texttt{config.json} 文件里
    \item 如果配置未设置好，智能体可以主动向用户询问所需信息
    \item 如果希望向用户展示\textbf{结构化的多选题}，可以让 Claude 使用 \texttt{AskUserQuestion} 工具
\end{itemize}

\subsection{description 字段是给模型看的（The Description Field Is For the Model）}

\begin{importantbox}{description 不是摘要，而是触发条件}
当 Claude Code 启动会话时，它会构建所有可用 Skills 及其 description 的清单。Claude 通过扫描这个清单来判断``这个请求有没有对应的 Skill？''这意味着 \texttt{description} 字段\textbf{不是功能摘要}（not a summary），而是\textbf{何时该触发这个 Skill 的描述}（a description of when to trigger）。好的 description 读起来更像一个 if-then 条件，而不是功能说明书。
\end{importantbox}

\subsection{本章小结}

编写 Skills 的核心原则：聚焦非显而易见的知识（避免重复 Claude 已知的内容）、持续积累 gotchas、利用文件夹结构做 progressive disclosure、避免过度约束（railroading）、设计好初始配置流程（config.json + AskUserQuestion）、把 description 当作触发条件而非功能说明。

